\chapter{Introduction}
\label{chap:introduction}

\section{Scope and terminology}
\label{sec:intro_scope}

The thesis title---\emph{AI for Control of Soft Robot by Using ML and Inverse Kinematics}---frames four keywords whose scope is fixed here. \emph{AI} and \emph{ML} are used interchangeably for learning-based models trained on recorded data, instantiated as four families: feed-forward regression (§\ref{sec:bg_ik_direct}), recurrent regression (§\ref{sec:bg_ik_temporal}), Jacobian-learned iteration (§\ref{sec:bg_ik_jac}), and sim-to-real pre-training (§\ref{sec:bg_sim}). \emph{Inverse Kinematics} denotes both the analytical piecewise-constant-curvature baseline (§\ref{sec:bg_models_pcc}) and the Jacobian-learned hybrid, which inverts an ML-trained forward map iteratively at run time. \emph{Control} is taken in its open-loop position-control sense---predict the actuation command that drives the tip to a target pose---following the usage of Giorelli et al.\ \cite{giorelli2015nn}, Ma et al.\ \cite{ma2022bp}, Fang et al.\ \cite{fang2022jacobian}, and Centurelli et al.\ \cite{centurelli2022closed}. Closed-loop real-hardware control is out of scope; the evaluation closes the loop in simulation through a shared learned forward-kinematics arbiter.

\section{Motivation}
\label{sec:intro_motivation}

Soft robots deform continuously under load rather than rotating rigid joints, which makes them safer for human interaction and confined spaces but hard to control: the actuation-to-pose map is nonlinear, history-dependent, and sensitive to manufacturing asymmetries of several centimetres \cite{manti2016soft}. Pneumatic soft manipulators---silicone cylinders with internal chambers pressurised by air---are the canonical instance and the focus here.

Four inverse-kinematics method families have been proposed independently over the last decade: analytical PCC inversion, feed-forward neural regression, recurrent neural regression, and Jacobian-learned iteration on a learned forward model \cite{webster2010review, giorelli2015nn, thuruthel2019modelbased, fang2022jacobian}. Each has been evaluated on the platform its authors introduced, on different metrics, so practitioners choosing between them must assemble non-comparable experiments across papers. This thesis consolidates the four families into a single experimental programme on two complementary pneumatic datasets under identical protocols.

\section{Problem, datasets, and evaluation}
\label{sec:intro_problem}

Given a target tip position $\mathbf{x}^{\ast} \in \mathbb{R}^{3}$, predict the actuation vector $\mathbf{p} \in \mathbb{R}^{n_{p}}$ that drives the tip there. Dataset~1 is a two-module I-Support recording \cite{manti2016soft, bianchi2023learning} with six chamber pressures mapping into a three-dimensional tip (multi-valued inverse); Dataset~2 is an undocumented three-valve recording with $n_{p} = 3$ (inverse at most single-valued on the observed workspace). Evaluation uses pressure-space prediction error, closed-loop position error through the FK arbiter, and mean tracking error on four smooth reference trajectories. No physical-hardware experiments are reported.

\section{Research questions and contributions}
\label{sec:intro_rqs}

\textbf{RQ1}: Does feed-forward neural IK outperform analytical PCC, and by how much? \textbf{RQ2}: Does adding a mid-marker triad reduce pressure error on the over-actuated arm in a statistically clean paired evaluation? \textbf{RQ3}: Does a recurrent architecture improve accuracy, and under what feature conditions? \textbf{RQ4}: Does Jacobian-learned iteration produce smaller closed-loop errors, and under what actuation redundancy? \textbf{RQ5}: Does a PyElastica-pre-trained network transfer to an undocumented platform?

Four contributions address the gaps in turn: (1)~a head-to-head comparison of the four IK families on two pneumatic datasets under identical protocols; (2)~a paired mid-marker ablation quantifying partial-shape sensing on identical test frames; (3)~a sim-to-real study at four real-data budgets with an explicit negative result attributable to workspace mismatch; (4)~a three-property decision criterion (actuation redundancy, sensing coverage, target-domain overlap) synthesising the regime-specific findings into a compact method-selection framework.

\section{Thesis structure}
\label{sec:intro_structure}

Chapter~\ref{chap:background} reviews the literature along the problem, modelling, and method axes; Chapter~\ref{chap:methods} describes datasets, methods, and protocols; Chapter~\ref{chap:results} reports the experimental programme; Chapter~\ref{chap:discussion} interprets the results regime by regime and states limitations and implications; Chapter~\ref{chap:conclusion} summarises, and Chapter~\ref{chap:future_work} identifies five concrete follow-ups.
